\chapter{Metode gradien proksimal}

% segment-id: d90.hab.v1.ch05.seg0001
Kita telah melihat dalam bukti (dan latihan praktis) bahwa metode subgradien tidak memberikan konvergensi yang ideal. Dalam kasus dasar atau terburuk, nilai fungsi berkonvergensi dengan orde \(\Oc(1/\sqrt{k})\), dan konvergensi memerlukan ukuran langkah yang mengecil. Dalam bab ini, kita akan memperkenalkan alternatif bagi subgradien (eksplisit) dalam optimisasi.

Perhatikan masalah tak berkendala
\begin{equation}\label{proximal:eq:problem}
    \min_x f(x),
\end{equation}
dengan \(f:\R^d\rightarrow(-\infty,\infty]\) proper, semikontinu bawah, dan konveks. Ingat pembaruan metode subgradien \(x_{k+1}=x_k-\tau_k g_k\), \(g_k\in\partial f(x_k)\), atau, secara ekuivalen, \(x_{k+1}\in x_k-\tau_k\partial f(x_k)\). Untuk \(\tau_k>0\), kita usulkan pembaruan implisit
\begin{equation}\label{proximal:eq:intro}
    x_{k+1}\in x_k-\tau_k\partial f(x_{k+1}).
\end{equation}
Pada tahap ini belum jelas apakah metode tersebut terdefinisi dengan baik. Namun, jika memang terdefinisi, secara intuitif kestabilannya dapat membaik karena arah pembaruan dipilih dengan \emph{melihat ke masa depan}. Gagasan ini selaras dengan kestabilan metode implisit untuk mendiskretkan persamaan diferensial (\cf~\emph{Euler implisit}).

Jika kita amati~\eqref{proximal:eq:intro} lebih dekat, aturan itu dapat ditulis sebagai
\begin{equation}
    0\in\partial\biggl(\frac{1}{2}\|\cdot-x_k\|^2+\tau_k f\biggr)(x_{k+1}),
\end{equation}
yang merupakan kondisi optimalitas untuk
\begin{equation}
    \min_x\frac{1}{2}\|x-x_k\|^2+\tau_k f(x).
\end{equation}
Hal ini membawa kita pada definisi berikut.
\begin{defn}[Operator proksimal]
    Operator atau pemetaan proksimal dari \(f:\R^d\rightarrow(-\infty,\infty]\) didefinisikan sebagai pemetaan
    \begin{equation}
        \begin{aligned}
            \prox_f:\R^d&\rightarrow 2^{\R^d},\\
            x&\mapsto\arg\min_{y\in\R^d}\frac{1}{2}\|y-x\|^2+f(y).
        \end{aligned}
    \end{equation}
\end{defn}

Kita biasanya menyebut \(\prox_f\) sebagai \enquote{proks}. Sebagai pemetaan bernilai himpunan, proks selalu dapat dituliskan; tetapi pemetaan itu baru berguna ketika \(\prox_f(x)\neq\emptyset\). Kita dapat membuktikan hasil berikut.

% segment-id: d90.hab.v1.ch05.seg0002
\begin{lemma}
    Misalkan \(f:\R^d\rightarrow(-\infty,\infty]\) proper, konveks, dan semikontinu bawah. Pemetaan proksimal bernilai tunggal. Lebih lanjut, \(y=\prox_f(x)\) jika dan hanya jika \(0\in(y-x)+\partial f(y)\).
\end{lemma}
\begin{rem}
    Berdasarkan karakterisasi \(0\in(y-x)+\partial f(y)\), kita sering menulis proks sebagai \(\prox_f(x)=(\Id+\partial f)^{-1}(x)\).
\end{rem}
\begin{proof}
    Demi kesederhanaan, andaikan \(\dom(f)^\circ\neq\emptyset\); untuk kasus umum, lihat~\cite{beck2017first}.

    Ambil \(g\in\partial f(y_0)\) untuk suatu \(y_0\in\dom(f)^\circ\). Kita mempunyai
    \begin{equation}
        f(y_0)+\inner{g}{y-y_0}+\frac{1}{2}\|x-y\|^2
        \leq f(y)+\frac{1}{2}\|x-y\|^2.
    \end{equation}
    Oleh karena itu (mengapa?), pemetaan \(y\mapsto f(y)+\frac{1}{2}\|x-y\|^2\) bersifat koersif. Pemetaan ini juga mewarisi sifat semikontinu bawah dan proper dari \(f\). Dengan metode langsung untuk fungsi proper, semikontinu bawah, dan koersif di ruang berdimensi hingga, terdapat solusi bagi
    \begin{equation}\label{proximal:eq:prox_proof}
        \min_{y\in\R^d}\frac{1}{2}\|y-x\|^2+f(y).
    \end{equation}
    Untuk keunikan, definisikan \(q_x(y)=\frac{1}{2}\|y-x\|^2+f(y)\) dan andaikan \(y_1\neq y_2\) adalah dua peminim. Karena suku kuadrat konveks ketat dan \(f\) konveks, untuk \(\bar y=(y_1+y_2)/2\) berlaku
    \begin{equation}
        \begin{aligned}
            q_x(\bar y)
            <&\frac{1}{2}\bigl(q_x(y_1)+q_x(y_2)\bigr)\\
            =&\min_{y\in\R^d}q_x(y),
        \end{aligned}
    \end{equation}
    sebuah kontradiksi. Pernyataan terakhir merupakan kondisi optimalitas bagi~\eqref{proximal:eq:prox_proof}.
\end{proof}

% segment-id: d90.hab.v1.ch05.seg0003
Setelah mengetahui bahwa proks terdefinisi dengan baik, kita dapat mendefinisikan metode gradien proksimal. Secara khusus, perhatikan masalah
\begin{equation}\label{proximal:eq:composite}
    \min_{x\in\R^d}F(x)\coloneqq f(x)+g(x),
\end{equation}
dengan \(f\) konveks dan \(L\)-halus, yakni terdiferensialkan dengan gradien \(L\)-Lipschitz, sedangkan \(g\) proper, konveks, dan semikontinu bawah. Masalah semacam ini sering muncul dalam praktik dan mencakup masalah semula dengan mengambil \(f\equiv0\). Untuk \(k=0,1,\dots\), metode gradien proksimal didefinisikan oleh
\begin{equation}\label{proximal:eq:prox_grad}
    \begin{aligned}
        x_{k+1}=\prox_{\tau_k g}\bigl(x_k-\tau_k\nabla f(x_k)\bigr).
    \end{aligned}
\end{equation}
Jadi, kita melakukan langkah gradien atau eksplisit terhadap \(f\), lalu langkah proksimal atau implisit terhadap \(g\).

Sebagai hasil sederhana pertama, kita menunjukkan bahwa titik tetap metode gradien proksimal tepat sama dengan titik optimal.
\begin{lemma}
    Untuk suatu \(\tau>0\) yang tetap dan \(x\in\R^d\), pernyataan berikut ekuivalen:
    \begin{equation}
        x=\prox_{\tau g}\bigl(x-\tau\nabla f(x)\bigr)
        \quad\Longleftrightarrow\quad
        0\in\partial(f+g)(x).
    \end{equation}
    Dengan demikian, \(x\) adalah titik tetap jika dan hanya jika \(x\) menyelesaikan~\eqref{proximal:eq:composite}.
\end{lemma}
\begin{proof}
    Karakterisasi optimalitas proks memberikan ekuivalensi
    \begin{equation}
        x=\prox_{\tau g}\bigl(x-\tau\nabla f(x)\bigr)
        \quad\Longleftrightarrow\quad
        0\in\tau\nabla f(x)+\tau\partial g(x).
    \end{equation}
    Karena \(\tau>0\), serta menurut aturan jumlah dan aturan skala untuk subdiferensial,
    \begin{equation}
        0\in\tau\bigl(\nabla f(x)+\partial g(x)\bigr)
        \quad\Longleftrightarrow\quad
        0\in\partial(f+g)(x).
    \end{equation}
    Untuk fungsi konveks \(F=f+g\), kondisi terakhir ekuivalen dengan optimalitas \(x\).
\end{proof}

% segment-id: d90.hab.v1.ch05.seg0004
Selain proks, kita mendefinisikan selubung Moreau sebagai nilai infimum dari objektif yang digunakan untuk proks.
\begin{defn}[Selubung Moreau]
    Misalkan \(f:\R^d\rightarrow(-\infty,\infty]\) dan \(\tau>0\). Selubung Moreau \(M_f^\tau:\R^d\rightarrow[-\infty,\infty]\) didefinisikan oleh
    \begin{equation}
        M_f^\tau(x)\coloneqq\inf_{y\in\R^d}\frac{1}{2\tau}\|x-y\|^2+f(y).
    \end{equation}
\end{defn}

\begin{lemma}
    Misalkan \(f:\R^d\rightarrow(-\infty,\infty]\) proper, konveks, dan semikontinu bawah, serta \(\tau>0\). Selubung Moreau selalu bernilai real (yakni berhingga), dan
    \begin{equation}
        M_f^\tau(x)
        =\frac{1}{2\tau}\bigl\|x-\prox_{\tau f}(x)\bigr\|^2
        +f\bigl(\prox_{\tau f}(x)\bigr).
    \end{equation}
\end{lemma}
\begin{proof}
    Latihan. Terapkan lemma keberadaan dan keunikan proks pada \(\tau f\), lalu substitusikan peminim unik tersebut ke dalam definisi infimum.
\end{proof}

\begin{example}
    Kita tinjau beberapa contoh pemetaan proksimal dan selubung Moreau yang sering muncul:
    \begin{enumerate}
        \item Misalkan \(C\subset\R^d\) tak kosong, tertutup, dan konveks, dan perhatikan fungsi indikator \(\delta_C\). Maka
        \begin{equation}
            \prox_{\tau\delta_C}(x)=\proj_C(x),
            \qquad
            M_{\delta_C}^{\tau}(x)=\frac{1}{2\tau}\operatorname{dist}(x,C)^2.
        \end{equation}
        \item Secara khusus, jika \(C=\{x\in\R^d\mid\|x\|_\infty\leq1\}\), maka
        \begin{equation}
            \bigl(\prox_{\tau\delta_C}(x)\bigr)_i
            =\frac{x_i}{\max\{1,|x_i|\}}.
        \end{equation}
        \item Jika \(f(x)=\|x\|_1\), pemetaan proksimalnya adalah operator \emph{ambang lunak}
        \begin{equation}
            \bigl(\prox_{\tau f}(x)\bigr)_i=
            \begin{cases}
                x_i-\tau & x_i>\tau,\\
                x_i+\tau & x_i<-\tau,\\
                0 & x_i\in[-\tau,\tau].
            \end{cases}
        \end{equation}
        Untuk \(d=1\), selubung Moreau-nya tepat merupakan fungsi Huber
        \begin{equation}
            M_{|\cdot|}^{\tau}(x)=
            \begin{cases}
                \dfrac{|x|^2}{2\tau} & |x|\leq\tau,\\
                |x|-\dfrac{\tau}{2} & |x|>\tau.
            \end{cases}
        \end{equation}
    \end{enumerate}
\end{example}

% segment-id: d90.hab.v1.ch05.seg0005
\begin{lemma}[Aturan komputasi untuk proks]
    Misalkan semua fungsi dasar yang muncul di bawah proper, semikontinu bawah, dan konveks pada domain berdimensi hingga yang dinyatakan. Aturan berikut berlaku:
    \begin{enumerate}
        \item Jika \(x=(x_1,\dots,x_m)\in\R^{d_1}\times\cdots\times\R^{d_m}\cong\R^d\), \(x_i\in\R^{d_i}\), \(\sum_i d_i=d\), dan \(f(x)=\sum_i f_i(x_i)\), maka \(\prox_f(x)=\bigl(\prox_{f_i}(x_i)\bigr)_i\).
        \item Jika \(f(x)=\alpha g(x)+b\) dengan \(\alpha>0\) dan \(b\in\R\), maka \(\prox_f=\prox_{\alpha g}\).
        \item Jika \(f(x)=g(\alpha x+b)\) dengan \(\alpha\neq0\) dan \(b\in\R^d\), maka \(\prox_f(x)=\alpha^{-1}\bigl(\prox_{\alpha^2g}(\alpha x+b)-b\bigr)\).
        \item Jika \(f(x)=g(Qx)\) dengan \(Q\in\R^{d\times d}\) ortogonal, maka \(\prox_f(x)=Q^\top\prox_g(Qx)\).
        \item Jika \(f(x)=g(x)+\inner{a}{x}+b\) dengan \(a\in\R^d\) dan \(b\in\R\), maka \(\prox_f(x)=\prox_g(x-a)\).
        \item Jika \(f(x)=g(x)+\frac{\gamma}{2}\|x-a\|^2\) dengan \(a\in\R^d\) dan \(\gamma\geq0\), maka \(\prox_f(x)=\prox_{\widetilde\gamma g}(\widetilde\gamma x+\gamma\widetilde\gamma a)\), dengan \(\widetilde\gamma=(1+\gamma)^{-1}\).
    \end{enumerate}
\end{lemma}

% segment-id: d90.hab.v1.ch05.seg0006
Selubung Moreau memberikan efek penghalusan pada fungsi \(f\), sebagaimana ditunjukkan oleh hasil berikut.
\begin{lemma}
    Misalkan \(f:\R^d\rightarrow(-\infty,\infty]\) proper, konveks, dan semikontinu bawah, serta \(\tau>0\). Selubung Moreau bersifat konveks dan terdiferensialkan. Gradiennya \(1/\tau\)-Lipschitz dan dapat ditulis sebagai
    \begin{equation}
        \nabla M_f^\tau(x)=\frac{1}{\tau}\bigl(x-\prox_{\tau f}(x)\bigr).
    \end{equation}
\end{lemma}
\begin{proof}
    Untuk \(q(x,y)=\frac{1}{2\tau}\|x-y\|^2+f(y)\), fungsi \(q\) konveks secara bersama pada \((x,y)\). Jika \(p_i=\prox_{\tau f}(x_i)\) dan \(t\in[0,1]\), evaluasi \(q\) di \(\bigl(tx_1+(1-t)x_2,tp_1+(1-t)p_2\bigr)\) menunjukkan bahwa \(M_f^\tau(tx_1+(1-t)x_2)\leq tM_f^\tau(x_1)+(1-t)M_f^\tau(x_2)\). Jadi, argumen ini tetap sah untuk fungsi \(f\) bernilai diperluas.

    Tetapkan \(p(x)=\prox_{\tau f}(x)\). Dari definisi proks,
    \begin{equation}
        \begin{aligned}
            M_f^\tau(x+h)
            =&\frac{1}{2\tau}\|x+h-p(x+h)\|^2+f(p(x+h))\\
            \leq&\frac{1}{2\tau}\|x+h-p(x)\|^2+f(p(x)),\\
            M_f^\tau(x)
            =&\frac{1}{2\tau}\|x-p(x)\|^2+f(p(x)).
        \end{aligned}
    \end{equation}
    Dengan mengurangkan dua hubungan itu, kita memperoleh
    \begin{equation}
        \begin{aligned}
            M_f^\tau(x+h)-M_f^\tau(x)
            \leq&\frac{1}{2\tau}\bigl(\|x+h-p(x)\|^2-\|x-p(x)\|^2\bigr)\\
            =&\inner{h}{\frac{1}{\tau}(x-p(x))}
              +\frac{1}{2\tau}\|h\|^2.
        \end{aligned}
    \end{equation}
    Secara ekuivalen,
    \begin{equation}
        \begin{aligned}
            M_f^\tau(x+h)-M_f^\tau(x)
            -\inner{h}{\frac{1}{\tau}(x-p(x))}
            \leq&\frac{1}{2\tau}\|h\|^2.
        \end{aligned}
    \end{equation}
    Definisikan residu keterdiferensialan
    \begin{equation}\label{proximal:eq:moreau_diff}
        r_x(h)\coloneqq
        M_f^\tau(x+h)-M_f^\tau(x)
        -\inner{h}{\frac{1}{\tau}(x-p(x))}.
    \end{equation}
    Kurangkan dua kondisi optimalitas proks, lalu gunakan monotonisitas \(\partial f\) dan ketaksamaan Cauchy--Schwarz. Hasilnya ialah \(\|p(x)-p(y)\|\leq\|x-y\|\), sehingga \(p\) tak ekspansif. Dengan mengevaluasi \(M_f^\tau(x)\) pada \(p(x+h)\), kita juga mendapatkan
    \begin{equation}
        r_x(h)
        \geq\inner{h}{\frac{1}{\tau}\bigl(p(x)-p(x+h)\bigr)}
             +\frac{1}{2\tau}\|h\|^2
        \geq-\frac{1}{2\tau}\|h\|^2.
    \end{equation}
    Secara keseluruhan,
    \begin{equation}
        \begin{aligned}
            \left|M_f^\tau(x+h)-M_f^\tau(x)
            -\inner{h}{\frac{1}{\tau}(x-p(x))}\right|
            \leq&\frac{1}{2\tau}\|h\|^2.
        \end{aligned}
    \end{equation}
    Jadi, \(\nabla M_f^\tau(x)=\tau^{-1}(x-p(x))\). Untuk membuktikan kekontinuan Lipschitz gradien, perhatikan bahwa
    \begin{equation}\label{proximal:eq:moreau_diff2}
        \begin{aligned}
            \|\nabla M_f^\tau(x)-\nabla M_f^\tau(y)\|^2
            =&\frac{1}{\tau}
              \inner{\nabla M_f^\tau(x)-\nabla M_f^\tau(y)}
                    {(x-p(x))-(y-p(y))}\\
            =&\frac{1}{\tau}
              \inner{\nabla M_f^\tau(x)-\nabla M_f^\tau(y)}{x-y}\\
             &-\frac{1}{\tau}
              \inner{\nabla M_f^\tau(x)-\nabla M_f^\tau(y)}{p(x)-p(y)}.
        \end{aligned}
    \end{equation}
    Kondisi optimalitas memberi
    \(\nabla M_f^\tau(x)=\tau^{-1}(x-p(x))\in\partial f(p(x))\).
    Karena itu, hasil kali dalam terakhir tidak negatif menurut monotonisitas subdiferensial, sehingga
    \begin{equation}\label{proximal:eq:moreau_diff2_bound}
        \begin{aligned}
            \|\nabla M_f^\tau(x)-\nabla M_f^\tau(y)\|^2
            \leq&\frac{1}{\tau}
              \inner{\nabla M_f^\tau(x)-\nabla M_f^\tau(y)}{x-y}\\
            \leq&\frac{1}{\tau}
              \|\nabla M_f^\tau(x)-\nabla M_f^\tau(y)\|\,\|x-y\|.
        \end{aligned}
    \end{equation}
    Setelah membagi dengan norma gradien yang berbeda dari nol (kasus nol langsung), diperoleh sifat \(1/\tau\)-Lipschitz.
\end{proof}

\begin{rem}
    Dengan representasi gradien selubung Moreau di atas, proks dapat ditulis kembali sebagai
    \begin{equation}
        \prox_{\tau f}(x)
        =x-\tau\biggl(\frac{1}{\tau}\bigl(x-\prox_{\tau f}(x)\bigr)\biggr)
        =x-\tau\nabla M_f^\tau(x).
    \end{equation}
    Jadi, satu langkah proksimal ekuivalen dengan satu langkah gradien pada selubung Moreau.
\end{rem}

\begin{example}
    Untuk fungsi indikator himpunan tak kosong, tertutup, dan konveks \(C\), rumusnya adalah \(\prox_{\tau\delta_C}(x)=\proj_C(x)\). Untuk \(f(x)=\|x\|_1\), setiap koordinat diberikan oleh ambang lunak \((\prox_{\tau f}(x))_i=\operatorname{sign}(x_i)\max\{|x_i|-\tau,0\}\), dengan \(\operatorname{sign}(0)=0\). Untuk \(f(x)=\|x\|_2\), penyusutan Euklides adalah \(\prox_{\tau f}(x)=(1-\tau/\|x\|_2)_+x\) jika \(x\neq0\), dan \(\prox_{\tau f}(0)=0\). Verifikasikan ketiga rumus ini langsung dari kondisi optimalitas proks.
\end{example}

% segment-id: d90.hab.v1.ch05.seg0007
Sekarang kita akan membuktikan konvergensi metode gradien proksimal. Mula-mula kita turunkan lemma penting: untuk fungsi \(L\)-halus, galat antara fungsi dan hampiran linearnya dibatasi dari atas oleh suku kuadrat.

\begin{lemma}\label{proximal_gradient:lemma:Lsmoothness}
    Misalkan \(F:\R^d\rightarrow\R\) terdiferensialkan kontinu dan gradiennya \(L\)-Lipschitz. Maka
    \begin{equation}
        F(y)\leq F(x)+\inner{\nabla F(x)}{y-x}+\frac{L}{2}\|y-x\|^2.
    \end{equation}
\end{lemma}
\begin{proof}
    Dari teorema dasar kalkulus dan sifat \(L\)-halus, kita mempunyai
    \begin{equation}
        \begin{aligned}
            F(y)-F(x)
            =&\int_0^1\frac{\dd}{\dd t}F(x+t(y-x))\dd t\\
            =&\int_0^1\inner{\nabla F(x+t(y-x))}{y-x}\dd t\\
            =&\int_0^1\inner{\nabla F(x+t(y-x))-\nabla F(x)}{y-x}\dd t
              +\int_0^1\inner{\nabla F(x)}{y-x}\dd t\\
            \leq&\int_0^1Lt\|y-x\|^2\dd t+\inner{\nabla F(x)}{y-x}\\
            =&\frac{L}{2}\|y-x\|^2+\inner{\nabla F(x)}{y-x}.
        \end{aligned}
    \end{equation}
\end{proof}

Untuk bukti konvergensi metode gradien proksimal, bagi \(\tau>0\) kita perkenalkan notasi
\begin{equation}
    T_\tau(x)=\frac{1}{\tau}\bigl(x-\prox_{\tau g}(x-\tau\nabla f(x))\bigr).
\end{equation}
Dengan fungsi ini, pembaruan gradien proksimal dapat ditulis sebagai
\begin{equation}
    x_{k+1}=x_k-\tau_kT_{\tau_k}(x_k).
\end{equation}

% segment-id: d90.hab.v1.ch05.seg0008
\begin{theorem}
    Misalkan \(f:\R^d\rightarrow\R\) konveks dan \(L\)-halus dengan \(L>0\), serta \(g:\R^d\rightarrow(-\infty,\infty]\) proper, semikontinu bawah, dan konveks. Andaikan \(\arg\min_xF(x)\neq\emptyset\) dan terdapat \(\tau_{\min}>0\) sedemikian sehingga
    \(0<\tau_{\min}\leq\tau_k\leq L^{-1}\) untuk semua \(k\). Untuk setiap \(x^*\in\arg\min_xF(x)\), algoritma gradien proksimal memenuhi
    \begin{equation}
        F(x_n)-F(x^*)\leq\frac{\|x_0-x^*\|^2}{2\tau_{\min}n}
        \qquad(n\geq1).
    \end{equation}
    Selain itu, \(x_k\) berkonvergensi menuju suatu \(\hat x\in\arg\min_xF(x)\).
\end{theorem}
\begin{proof}
    Dari definisi \(T_\tau\), kita mempunyai
    \begin{equation}
        T_{\tau_k}(x_k)-\nabla f(x_k)
        \in\partial g\bigl(x_k-\tau_kT_{\tau_k}(x_k)\bigr)
        =\partial g(x_{k+1}).
    \end{equation}
    Dengan memakai hubungan ini, kekonveksan \(f\), Lemma~\ref{proximal_gradient:lemma:Lsmoothness}, dan \(\tau_k\leq L^{-1}\), untuk setiap \(z\in\R^d\) diperoleh
    \begin{equation}
        \begin{aligned}
            F(x_{k+1})
            \leq& f(x_k)-\tau_k\inner{\nabla f(x_k)}{T_{\tau_k}(x_k)}
              +\frac{\tau_k}{2}\|T_{\tau_k}(x_k)\|^2\\
             &+g(z)-\inner{T_{\tau_k}(x_k)-\nabla f(x_k)}
                               {z-\bigl(x_k-\tau_kT_{\tau_k}(x_k)\bigr)}\\
            \leq& f(z)+g(z)-\inner{\nabla f(x_k)}{z-x_k}
              -\tau_k\inner{\nabla f(x_k)}{T_{\tau_k}(x_k)}
              +\frac{\tau_k}{2}\|T_{\tau_k}(x_k)\|^2\\
             &-\inner{T_{\tau_k}(x_k)-\nabla f(x_k)}
                               {z-\bigl(x_k-\tau_kT_{\tau_k}(x_k)\bigr)}\\
            =&F(z)-\inner{T_{\tau_k}(x_k)}{z-x_k}
              -\frac{\tau_k}{2}\|T_{\tau_k}(x_k)\|^2.
        \end{aligned}
    \end{equation}
    Dengan mengambil \(z=x_k\), kita memperoleh
    \begin{equation}
        F(x_{k+1})\leq F(x_k)
        -\frac{\tau_k}{2}\|T_{\tau_k}(x_k)\|^2.
    \end{equation}
    Jadi, nilai objektif tidak meningkat dan turun ketat kecuali jika \(T_{\tau_k}(x_k)=0\); menurut lemma titik tetap, kasus terakhir berarti \(x_k\) optimal. Di sisi lain, ambil \(z=x^*\) dan gunakan \(T_{\tau_k}(x_k)=\tau_k^{-1}(x_k-x_{k+1})\). Maka
    \begin{equation}
        \begin{aligned}
            0\leq F(x_{k+1})-F(x^*)
            \leq&\inner{T_{\tau_k}(x_k)}{x_k-x^*}
                 -\frac{\tau_k}{2}\|T_{\tau_k}(x_k)\|^2\\
            =&\frac{1}{2\tau_k}
              \left(2\inner{x_k-x_{k+1}}{x_k-x^*}
                    -\|x_k-x_{k+1}\|^2\right)\\
            =&\frac{1}{2\tau_k}
              \left(\|x_k-x^*\|^2-\|x_{k+1}-x^*\|^2\right).
        \end{aligned}
    \end{equation}
    Khususnya, \(\|x_k-x^*\|^2-\|x_{k+1}-x^*\|^2\geq0\), sehingga jarak ke setiap peminim tidak meningkat. Dengan menjumlahkan untuk \(k=0,\dots,n-1\), kita memperoleh
    \begin{equation}
        \begin{aligned}
            \sum_{k=0}^{n-1}\bigl(F(x_{k+1})-F(x^*)\bigr)
            \leq&\sum_{k=0}^{n-1}\frac{1}{2\tau_k}
              \left(\|x_k-x^*\|^2-\|x_{k+1}-x^*\|^2\right)\\
            \leq&\frac{1}{2\tau_{\min}}\sum_{k=0}^{n-1}
              \left(\|x_k-x^*\|^2-\|x_{k+1}-x^*\|^2\right)\\
            =&\frac{1}{2\tau_{\min}}
              \left(\|x_0-x^*\|^2-\|x_n-x^*\|^2\right).
        \end{aligned}
    \end{equation}
    Karena \(F(x_k)-F(x^*)\) juga tidak meningkat,
    \begin{equation}
        n\bigl(F(x_n)-F(x^*)\bigr)
        \leq\sum_{k=0}^{n-1}\bigl(F(x_{k+1})-F(x^*)\bigr)
        \leq\frac{\|x_0-x^*\|^2}{2\tau_{\min}}.
    \end{equation}
    Jadi,
    \begin{equation}
        F(x_n)-F(x^*)\leq
        \frac{\|x_0-x^*\|^2}{2\tau_{\min}n}.
    \end{equation}
    Terakhir, kita buktikan konvergensi \((x_k)_k\). Perhitungan di atas menunjukkan bahwa \(\|x_k-x^*\|\) tidak meningkat untuk setiap peminim \(x^*\). Dengan demikian, \((x_k)_k\) terbatas dan, di ruang berdimensi hingga, mempunyai subbarisan konvergen.
    Misalkan \(\hat x=\lim_jx_{n(j)}\) adalah suatu titik akumulasi. Karena \(F(x_k)\rightarrow\min F\) dan \(F=f+g\) semikontinu bawah,
    \[
        F(\hat x)\leq\liminf_jF(x_{n(j)})=\lim_kF(x_k)=\min F.
    \]
    Jadi, \(\hat x\) adalah peminim dan \(\|x_k-\hat x\|\) tidak meningkat. Karena itu, terdapat \(c\geq0\) dengan \(\|x_k-\hat x\|\rightarrow c\). Namun,
    \begin{equation}
        c=\lim_k\|x_k-\hat x\|
         =\lim_j\|x_{n(j)}-\hat x\|=0.
    \end{equation}
    Maka seluruh barisan berkonvergensi ke \(\hat x\), dan bukti selesai.
\end{proof}
